\documentclass{article}
\usepackage[utf8]{inputenc}
\usepackage[spanish]{babel}
\usepackage{amsmath, amssymb}
\usepackage{graphicx}
\usepackage{hyperref}

\newcommand{\R}{\mathbb{R}}
\newtheorem{teorema}{Teorema}

\title{Mars: LaTeX sin compilar}
\author{Tu nombre}
\date{\today}

\begin{document}

\maketitle

\begin{abstract}
Este documento muestra lo que el editor visual sabe dibujar. Haz clic en cualquier texto y
escribe encima, como en un procesador de textos: el archivo \texttt{.tex} se actualiza solo.
\end{abstract}

\tableofcontents

\section{Texto} \label{sec:texto}

Un párrafo con \textbf{negrita}, \textit{cursiva}, \underline{subrayado} y \texttt{código}.
Los acentos funcionan: canción, niño, ¿qué tal? % este comentario no se pierde al editar
También las comillas ``dobles'', los guiones -- y --- y los símbolos 50\% y \$10.

Una fórmula en línea como $e^{i\pi} + 1 = 0$ se edita con doble clic. Mira la
sección~\ref{sec:mates} y la ecuación~\eqref{eq:gauss}\footnote{Las notas al pie se juntan al final.}.

\subsection{Listas}

\begin{itemize}
  \item Primer elemento
  \item Segundo elemento con $x \in \R$
  \begin{enumerate}
    \item Anidado uno
    \item Anidado dos
  \end{enumerate}
\end{itemize}

\section{Matemáticas} \label{sec:mates}

\begin{equation} \label{eq:gauss}
  \int_{-\infty}^{\infty} e^{-x^2} \, dx = \sqrt{\pi}
\end{equation}

\begin{align}
  (a+b)^2 &= a^2 + 2ab + b^2 \\
  (a-b)^2 &= a^2 - 2ab + b^2
\end{align}

\begin{teorema}[Pitágoras]
En un triángulo rectángulo, $a^2 + b^2 = c^2$.
\end{teorema}

\section{Tablas y código}

\begin{table}[h]
  \centering
  \begin{tabular}{|l|c|r|}
    \hline
    \textbf{Planeta} & \textbf{Lunas} & \textbf{Radio (km)} \\
    \hline
    Tierra & 1 & 6371 \\
    \hline
    Marte & 2 & 3390 \\
    \hline
  \end{tabular}
  \caption{Dos planetas rocosos}
  \label{tab:planetas}
\end{table}

\begin{verbatim}
def hola():
    print("Hola, Marte")
\end{verbatim}

\end{document}
